\section{Background}\label{sec:background}

\subsection{Presentations and their semantics}\label{sec:presentations}

A \emph{relation} on an alphabet $X$ is a set of pairs of words over $X$.
It generates a congruence $\approx$ on the free monoid $X^*$: the
smallest equivalence relation that contains the pairs and is closed under
concatenation on both sides. A \emph{semantics} is a monoid homomorphism
$\llbracket - \rrbracket : X^* \to M$. It is \emph{sound} if related
words have equal images, and \emph{complete} if words with equal images
are related. When it is sound and complete, the relation presents the
image of $\llbracket - \rrbracket$; when moreover every element of $M$
is the image of a word, it presents $M$. All the presentation theorems in this
paper have one of these two forms. For circuits, $M$ is a group of matrices
and the image is the group the gates generate, which in most cases has
no simpler description. For the two-level groups, $M$ is a matrix group
such as $O_n(\Z[\frac1{\sqrt2}])$ and the generators generate all of it.

Our library~\cite{anonymous2027popl} fixes this vocabulary, and the
syntax of circuits (Section~\ref{sec:library}). A circuit on $n$
wires is a word over wire-indexed gates. A gate of arity $k$ sits on the
bottom $k$ wires, and the shift $\uparrow$ moves a circuit up one wire.
The library extends any gate relation with four structural rules: an
equation holds one wire up (\aid{cong↑}), a one- or two-wire gate at the
bottom commutes with anything shifted above it (\aid{comm₁},
\aid{comm₂}), and a scalar is the same on every wire (\aid{ω↑=ω}). That
a scalar commutes with every gate is not a rule. It follows from the
four when there is one scalar, and we show in
Section~\ref{sec:structural} that none of the four follows from the
other three.

Many of the papers we formalized present groups of \emph{two-level}
matrices. Over a ring $R$ with a suitable element $s$ (for example
$s=\frac1{\sqrt2}$), the generators act on $R^n$: $(-1)_{[a]}$ negates
coordinate~$a$, $X_{[a,b]}$ exchanges coordinates $a$ and $b$, and
$H_{[a,b]}$ applies $s\left(\begin{smallmatrix}1&1\\1&-1\end{smallmatrix}\right)$
to coordinates $a$ and~$b$. Circuit groups embed into such groups, which
is how several circuit completeness proofs reduce to them.

\subsection{Completeness by normal forms}\label{sec:by-nf}

The first method defines a set of normal forms with an interpretation
as words, shows that every word is related to the interpretation of
some normal form, and shows that the semantics is injective on normal
forms, so that two words with the same matrix are related to the same
normal form. The second step is the long one: the usual proof pushes
one gate at a time into a normal form, and needs a rewrite rule for
every gate and every shape of normal form, the ``typed relations'' or
tables that fill the appendices. The third step is usually semantic:
one reads the normal form back from the matrix, for example from how
it acts on Pauli operators.

\subsection{Completeness by exact synthesis and levels}\label{sec:by-levels}

The second method starts from an exact synthesis algorithm in the style
of Giles and Selinger~\cite{giles2013exact}. Let $M$ be a matrix in the
group. Its \emph{pivot} is the last column that differs from the
corresponding column of the identity, and the \emph{least denominator
exponent} $k$ of that column is the least $k$ such that $\sqrt2^k$ times
the column has entries in $\Z[\sqrt2]$, or in the integral ring of
the case at hand. The algorithm looks at the residues of these
entries and chooses a short word, the \emph{syllable}, that lowers a
\emph{level}: in Hadamard-$\Pi$ the level is the triple of the pivot,
$k$, and the number of odd entries, ordered lexicographically. Iterating
gives a word $\mathit{nw}(M)$ with $\mathit{nw}(M)\cdot M = I$. This is
generation: every matrix is the matrix of a word.

For completeness one shows that, for every word $w$ and matrix $M$,
\[
  \mathit{nw}(w\cdot M) \, w \;\approx\; \mathit{nw}(M).
\]
Call such a $w$ a \emph{path} from $M$. Applied to $M = I$ this gives
$\mathit{nw}(\llbracket w \rrbracket)\, w \approx \varepsilon$, so two
words with the same matrix are related to the same word, after
cancelling $\mathit{nw}(\llbracket w \rrbracket)$. Paths compose, so it suffices to show that every \emph{edge}
$M \to g \cdot M$ is a path. The proof is by well-founded induction on
the level of the edge, the larger of the levels of its two ends. If $g$
is the syllable of $M$, or of $g\cdot M$, the edge is a path by
definition. Otherwise one completes a square: with $N$ the syllable of
$M$ and $N'$ that of $g\cdot M$, one finds a word $W$ with
$W N' g \approx N$ whose intermediate states all lie below the level of
the edge, so that by induction every edge of $W$ is a path, and the
square closes. The \emph{Main Lemma} of these papers is the statement
that such squares exist for all states and generators. Its proof is
the long case analysis, each case a diagram of matrices whose levels
must be computed, and most of the effort goes into it. The rest,
generation, soundness and the reduction to edges, is short and the
same each time.
